\subsection{Protocol-dependent gaps and initial-state dependence}
\label{sec:rewrite-gaps}

For an invariant space $\mathcal L$, let $P_0^{\mathcal L}(s)$ be the
ground projector of $H(s)|_{\mathcal L}$ and
$g_0^{\mathcal L}(s)=\operatorname{rank}P_0^{\mathcal L}(s)$.  Write the
eigenvalues of the restriction, ordered with multiplicity, as
$E_0^{\mathcal L}(s)\leq E_1^{\mathcal L}(s)\leq\cdots$.  Whenever the
indicated complementary level exists, we keep separate
\begin{align}
 \delta_{\rm mult}^{\mathcal L}(s)
 &=E_1^{\mathcal L}(s)-E_0^{\mathcal L}(s),\\
 \Delta_{\rm exc}^{\mathcal L}(s)
 &=E_{g_0^{\mathcal L}(s)}^{\mathcal L}(s)-E_0^{\mathcal L}(s).
 \label{eq:two-gaps-rewrite}
\end{align}
If $\dim\mathcal L=1$, we set $\delta_{\rm mult}^{\mathcal L}=+\infty$;
if $g_0^{\mathcal L}=\dim\mathcal L$, we set
$\Delta_{\rm exc}^{\mathcal L}=+\infty$.  These conventions express that
the corresponding complementary level does not exist.
The first vanishes whenever the ground state is degenerate; the second is
the excitation gap above the whole manifold.  A Schur multiplicity or an
exact coincidence between irreps is therefore not called a gap closure when
Eq.~\eqref{eq:two-gaps-rewrite}'s second quantity stays positive.

Let $E_0$, $E_0^{\rm sym}$, and $E_0^{\mathcal K}$ be the ground energies
of $H(s)$ on $\mathcal H_{\rm addr}$, $\mathcal H_{\rm sym}$, and
$\mathcal K_u$, respectively, at a common value of $s$.  The nested
invariant spaces in Eq.~\eqref{eq:nested-spaces-rewrite} define the
restriction costs
\begin{align}
 \delta_{\rm sym}&=E_0^{\rm sym}-E_0,\quad\quad
 \delta_{\rm cyc}=E_0^{\mathcal K}-E_0^{\rm sym},\\
 \delta_{\rm acc}&=E_0^{\mathcal K}-E_0
 =\delta_{\rm sym}+\delta_{\rm cyc}.
 \label{eq:accessibility-cost-rewrite}
\end{align}
A positive $\delta_{\rm acc}$ proves that the common cyclic envelope excludes
the global ground branch.  If an isolated band is then followed inside that
envelope, it is globally excited but can nevertheless have high cover
probability.  Conversely, zero restriction cost does not prove that a
particular schedule prepares the corresponding ground vector.

For adiabatic evolution, the relevant object is an isolated spectral band,
not merely the first distinct final eigenvalue
\cite{JansenRuskaiSeiler2007}.  Let
$P_{\rm track}(s)$ be a continuous isolated spectral projector on
$\mathcal K_u$ whose range contains the initial ground state.  Set
$H_{\mathcal K_u}(s)=H(s)|_{\mathcal K_u}$ and let
$H_{\rm track}(s)$ be the compression of $H_{\mathcal K_u}(s)$ to
$\Ran P_{\rm track}(s)$.  Its complementary compression is
\begin{equation}
 H_\perp(s)=H_{\mathcal K_u}(s)
 \big|_{\Ran(I_{\mathcal K_u}-P_{\rm track}(s))},
\end{equation}
where $I_{\mathcal K_u}$ is the identity on $\mathcal K_u$.  For finite
real sets $A,B$, write
$\operatorname{dist}(A,B)=\min_{a\in A,b\in B}|a-b|$ and adopt the
convention $\operatorname{dist}(A,\varnothing)=+\infty$.  Define
\begin{multline}
 \Delta_{\rm ad}(s)=\operatorname{dist}\!\left(
 \spec H_{\rm track}(s),\right.\\
 \left.\spec H_\perp(s)\right).
 \label{eq:adiabatic-gap-rewrite}
\end{multline}
This projector-based definition retains multiplicity: if a complementary
eigenvector has the same numerical eigenvalue as the tracked band, then
$\Delta_{\rm ad}=0$.
An isolated continuous projector has constant rank.  If its intended rank
changes, an isolation gap has closed and a nondegenerate adiabatic formula
cannot simply be continued through that point.

The four quantities have the following physical meanings:
\begin{center}
\centering
\small
\begin{tabular}{lp{0.62\columnwidth}}
\hline
quantity & physical meaning\\
\hline
$\delta_{\rm mult}$ & splitting inside the lowest manifold; zero for a
degenerate ground energy\\
$\Delta_{\rm exc}$ & excitation gap above the entire lowest manifold\\
$\delta_{\rm acc}$ & energy cost of confinement to the common cyclic
envelope\\
$\Delta_{\rm ad}$ & isolation of the spectral band followed inside that
space\\
\hline
\end{tabular}
\end{center}

Both the envelope and the tracked band are conditional on the protocol.  If
$[H(t),\Pi_\gamma]=0$ and the initial state lies in $\Ran\Pi_\gamma$, its
evolution remains there, but an initial state adapted to another irrep or a
symmetry-breaking generator changes the invariant and tracked spectra.  This statement
does not forbid transient concentration in a computationally marked
subspace: it forbids only transition amplitude into symmetry sectors
orthogonal to the one selected by the initial state.
